% Options for packages loaded elsewhere
\PassOptionsToPackage{unicode}{hyperref}
\PassOptionsToPackage{hyphens}{url}
%
\documentclass[
  14pt,
]{extarticle}
\usepackage{lmodern}
\usepackage{amssymb,amsmath}
\usepackage{ifxetex,ifluatex}
\ifnum 0\ifxetex 1\fi\ifluatex 1\fi=0 % if pdftex
  \usepackage[T1]{fontenc}
  \usepackage[utf8]{inputenc}
  \usepackage{textcomp} % provide euro and other symbols
\else % if luatex or xetex
  \usepackage{unicode-math}
  \defaultfontfeatures{Scale=MatchLowercase}
  \defaultfontfeatures[\rmfamily]{Ligatures=TeX,Scale=1}
  \setmainfont[]{Droid Sans Fallback}
\fi
% Use upquote if available, for straight quotes in verbatim environments
\IfFileExists{upquote.sty}{\usepackage{upquote}}{}
\IfFileExists{microtype.sty}{% use microtype if available
  \usepackage[]{microtype}
  \UseMicrotypeSet[protrusion]{basicmath} % disable protrusion for tt fonts
}{}
\makeatletter
\@ifundefined{KOMAClassName}{% if non-KOMA class
  \IfFileExists{parskip.sty}{%
    \usepackage{parskip}
  }{% else
    \setlength{\parindent}{0pt}
    \setlength{\parskip}{6pt plus 2pt minus 1pt}}
}{% if KOMA class
  \KOMAoptions{parskip=half}}
\makeatother
\usepackage{xcolor}
\IfFileExists{xurl.sty}{\usepackage{xurl}}{} % add URL line breaks if available
\IfFileExists{bookmark.sty}{\usepackage{bookmark}}{\usepackage{hyperref}}
\hypersetup{
  hidelinks,
  pdfcreator={LaTeX via pandoc}}
\urlstyle{same} % disable monospaced font for URLs
\usepackage[margin=1in]{geometry}
\usepackage{longtable,booktabs}
% Correct order of tables after \paragraph or \subparagraph
\usepackage{etoolbox}
\makeatletter
\patchcmd\longtable{\par}{\if@noskipsec\mbox{}\fi\par}{}{}
\makeatother
% Allow footnotes in longtable head/foot
\IfFileExists{footnotehyper.sty}{\usepackage{footnotehyper}}{\usepackage{footnote}}
\makesavenoteenv{longtable}
\setlength{\emergencystretch}{3em} % prevent overfull lines
\providecommand{\tightlist}{%
  \setlength{\itemsep}{0pt}\setlength{\parskip}{0pt}}
\setcounter{secnumdepth}{-\maxdimen} % remove section numbering

\author{}
\date{}

\begin{document}

\hypertarget{week-07-ux4f5cux696dux56deux994b}{%
\section{Week 07 作業回饋}\label{week-07-ux4f5cux696dux56deux994b}}

\begin{quote}
\textbf{評分標準：} 依據 \texttt{week07\_rubric.xlsx}
之六項評分指標，基礎分 100 分，Bonus 最高 +10 分，總分上限為 100。
各大項配分：Q1 Data Simulation (25) ／ Q2 Data Cleaning (20) ／ Q3
Descriptive Stats (20) ／ Q4 Multi-Panel Figure (35) ／ Bonus Panel E
(+5) ／ Bonus Paired t-test (+5)
\end{quote}

\begin{center}\rule{0.5\linewidth}{0.5pt}\end{center}

\hypertarget{ux5442ux6770ux9a5b-100100}{%
\subsection{113825002 呂杰驛 ---
100／100}\label{ux5442ux6770ux9a5b-100100}}

\hypertarget{ux7e3dux8a55}{%
\subsubsection{總評}\label{ux7e3dux8a55}}

完整度為本次作業最高的一位同學之一：四道基礎題與兩項 Bonus（Panel E
練習/疲勞曲線、paired-samples t-test）皆完成，且在清理流程、pivot\_table
計算、以及 5 個子圖的版面均表現穩定。因兩項 Bonus
抵消了基礎題的若干規格偏差（詳見下表），總分達到上限
100。建議下次在圖表細節（圖總標題、Panel B
文字註解）與欄位命名一致性上再做精修。

\hypertarget{ux5404ux9805ux8a55ux8a9e}{%
\subsubsection{各項評語}\label{ux5404ux9805ux8a55ux8a9e}}

\begin{longtable}[]{@{}lll@{}}
\toprule
\begin{minipage}[b]{0.10\columnwidth}\raggedright
指標\strut
\end{minipage} & \begin{minipage}[b]{0.04\columnwidth}\raggedright
得分（Base）\strut
\end{minipage} & \begin{minipage}[b]{0.77\columnwidth}\raggedright
評語\strut
\end{minipage}\tabularnewline
\midrule
\endhead
\begin{minipage}[t]{0.10\columnwidth}\raggedright
Q1 Data Simulation (NumPy)\strut
\end{minipage} & \begin{minipage}[t]{0.04\columnwidth}\raggedright
24/25\strut
\end{minipage} & \begin{minipage}[t]{0.77\columnwidth}\raggedright
\texttt{np.random.default\_rng(2026)} ✓，20 人 × 2 條件 × 40 trials
✓，分配參數 Normal(520,80)／Normal(620,100)、Binomial(1, 0.95)／(1,
0.85) 皆正確。唯受試者欄位命名為 \texttt{ID} 而非題目規範的
\texttt{subject}，此命名差異在後續 groupby／pivot\_table
都一致使用，不影響結果，但對協作與 CSV 交換會造成格式不一致，扣 1
分。\strut
\end{minipage}\tabularnewline
\begin{minipage}[t]{0.10\columnwidth}\raggedright
Q2 Data Cleaning\strut
\end{minipage} & \begin{minipage}[t]{0.04\columnwidth}\raggedright
20/20\strut
\end{minipage} & \begin{minipage}[t]{0.77\columnwidth}\raggedright
三步驟依序執行（errors → per-subject outliers → implausible
RTs），per-subject 離群值以
\texttt{groupby("ID"){[}"rt\_ms"{]}.transform("mean"/"std")}
正確實作，清理摘要完整印出各步驟移除數、最終數與保留百分比，並存出
\texttt{stroop\_clean.csv}（\texttt{index=False} ✓）。\strut
\end{minipage}\tabularnewline
\begin{minipage}[t]{0.10\columnwidth}\raggedright
Q3 Descriptive Stats (groupby)\strut
\end{minipage} & \begin{minipage}[t]{0.04\columnwidth}\raggedright
20/20\strut
\end{minipage} & \begin{minipage}[t]{0.77\columnwidth}\raggedright
使用
\texttt{agg({[}\textquotesingle{}mean\textquotesingle{},\textquotesingle{}std\textquotesingle{},\textquotesingle{}median\textquotesingle{},\textquotesingle{}count\textquotesingle{}{]})}
計算每受試者 × 條件統計；以 \texttt{pivot\_table} 重塑後做
\texttt{incongruent\ −\ congruent} 得到 interference
cost；\texttt{scipy.stats.sem} 計算 between-subjects SEM，SEM／SD／Range
之解讀皆正確。\strut
\end{minipage}\tabularnewline
\begin{minipage}[t]{0.10\columnwidth}\raggedright
Q4 Multi-Panel Figure\strut
\end{minipage} & \begin{minipage}[t]{0.04\columnwidth}\raggedright
31/35\strut
\end{minipage} & \begin{minipage}[t]{0.77\columnwidth}\raggedright
四個子圖皆完成且資訊量充足：Panel A histogram + 平均線、Panel C 橫條圖依
cost 由大至小排序且畫出組平均虛線、Panel D
依條件著色散佈圖。兩項可再改進：（1）圖總標題為
\texttt{"Summary\ for\ Stroop\ Task"}，題目規範為
\texttt{"Stroop\ Task:\ Experiment\ Summary"}（−2）；（2）Panel B
的文字註解顯示的是 SEM
值（\texttt{f\textquotesingle{}\{sems{[}i{]}:.2f\}\textquotesingle{}}），題目要求標注「mean
value above each bar」，建議改為
\texttt{f\textquotesingle{}\{bar.get\_height():.1f\}\textquotesingle{}}（−2）。\texttt{tight\_layout(rect={[}0,0,1,0.96{]})}
與 PNG/PDF 雙儲存皆正確。\strut
\end{minipage}\tabularnewline
\begin{minipage}[t]{0.10\columnwidth}\raggedright
\textbf{Base 小計}\strut
\end{minipage} & \begin{minipage}[t]{0.04\columnwidth}\raggedright
\textbf{95/100}\strut
\end{minipage} & \begin{minipage}[t]{0.77\columnwidth}\raggedright
\strut
\end{minipage}\tabularnewline
\begin{minipage}[t]{0.10\columnwidth}\raggedright
Bonus Panel E (+5)\strut
\end{minipage} & \begin{minipage}[t]{0.04\columnwidth}\raggedright
\textbf{+4}\strut
\end{minipage} & \begin{minipage}[t]{0.77\columnwidth}\raggedright
Panel E 線圖正確呈現 congruent／incongruent 兩條件逐 trial 的平均
RT，\texttt{grid} 與 \texttt{legend} 齊全。唯標題中
\texttt{"\textbackslash{}n"} 誤寫為
\texttt{"/n"}（\texttt{\textquotesingle{}(E)\ Mean\ RT\ Across\ Trial\ Positions\ /nin\ Different\ Conditions\textquotesingle{}}），實際輸出變成字面上的
\texttt{/nin}，扣 1 分。\strut
\end{minipage}\tabularnewline
\begin{minipage}[t]{0.10\columnwidth}\raggedright
Bonus Paired t-test (+5)\strut
\end{minipage} & \begin{minipage}[t]{0.04\columnwidth}\raggedright
\textbf{+5}\strut
\end{minipage} & \begin{minipage}[t]{0.77\columnwidth}\raggedright
正確以 \texttt{scipy.stats.ttest\_rel(incongruent,\ congruent)}
執行配對樣本 t 檢定，輸出 t 值與 p 值並附 α=.05
的口語判讀，統計方向解釋無誤。\strut
\end{minipage}\tabularnewline
\begin{minipage}[t]{0.10\columnwidth}\raggedright
\textbf{最終分數（上限 100）}\strut
\end{minipage} & \begin{minipage}[t]{0.04\columnwidth}\raggedright
\textbf{100/100}\strut
\end{minipage} & \begin{minipage}[t]{0.77\columnwidth}\raggedright
Base 95 + Bonus 9 = 104 → 封頂 100。\strut
\end{minipage}\tabularnewline
\bottomrule
\end{longtable}

\hypertarget{ux5f8cux7e8cux5efaux8b70}{%
\subsubsection{後續建議}\label{ux5f8cux7e8cux5efaux8b70}}

\begin{enumerate}
\def\labelenumi{\arabic{enumi}.}
\tightlist
\item
  \textbf{欄位命名一致性}：資料檔 schema 以規範
  \texttt{subject}、\texttt{trial}、\texttt{condition}、\texttt{rt\_ms}、\texttt{correct}
  為原則，自行改名會使 CSV 難以被其他腳本（或共同作者）直接重用。
\item
  \textbf{圖表標題規範}：Matplotlib 中換行字元是反斜線
  \texttt{\textbackslash{}n}，不是斜線 \texttt{/n}；另外
  \texttt{suptitle} 與題目要求的標題文字建議完全一致以利批改比對。
\item
  \textbf{圖表可讀性}：Panel B
  在同時顯示誤差棒與文字註解時，建議文字註解顯示的是「主要測量量」（mean），誤差已經由
  errorbar 視覺化，另行標示 SEM 數值反而容易讓讀者誤會長條高度的意義。
\end{enumerate}

\end{document}
